\section{Urgency: patients are waiting now}
\label{sec:urgency}

The final appeal is timing. Legislation accelerates when delay is shown to have a
cost: patients are waiting now, lives are being lost now, and the technology is
advancing faster than the rules that govern it. Urgency is only persuasive when it
is paired with a credible remedy, so that a committee feels able to act rather
than overwhelmed.

[DRAFTING INSTRUCTIONS: in the full-narrative, anchor urgency in
\texttt{review/references/narrative\_refs.bib} entries \texttt{kok2013finding}
(fear appeals work only when paired with self-efficacy, so always present the
remedy) and \texttt{fda2021samd}. Reproduce Mermaid figures
\texttt{review/mermaid/diagrams/16-urgency-capability-gap.md} and
\texttt{21-national-platform-capstone.md} as colored TikZ. Close the section, and
the paper, with the single call to action: enact H. R. 9510 v5.0 (DOI
10.5281/zenodo.20619762). Add a body-width table of the cost of each year of
delay against the cost of acting.]

\begin{narrativefig}
\node[nftwo] (a) at (0,0) {Capability rising};
\node[nfrisk] (b) at (3.8,0) {Regulation gap};
\node[nfact] (c) at (7.6,0) {H. R. 9510};
\node[nfhope] (d) at (10.8,0) {Act now};
\draw[nfedge] (a)--(c); \draw[nfedge] (b)--(c); \draw[nfedge] (c)--(d);
\end{narrativefig}
\vspace{-0.2cm}
\figcaption{Figure 9. Closing the gap before more time is lost (placeholder;
expanded from Mermaid 16 in the full-narrative).}

The eight appeals converge on one request. The evidence exists, the tools are
tested, and the people are waiting. The committee has the power, and now the
reason, to act.
